% Options for packages loaded elsewhere
\PassOptionsToPackage{unicode}{hyperref}
\PassOptionsToPackage{hyphens}{url}
\PassOptionsToPackage{dvipsnames,svgnames,x11names}{xcolor}
%
\documentclass[
  10pt,
  a4paper,
]{article}
\usepackage{amsmath,amssymb}
\usepackage{iftex}
\ifPDFTeX
  \usepackage[T1]{fontenc}
  \usepackage[utf8]{inputenc}
  \usepackage{textcomp} % provide euro and other symbols
\else % if luatex or xetex
  \usepackage{unicode-math} % this also loads fontspec
  \defaultfontfeatures{Scale=MatchLowercase}
  \defaultfontfeatures[\rmfamily]{Ligatures=TeX,Scale=1}
\fi
\usepackage{lmodern}
\ifPDFTeX\else
  % xetex/luatex font selection
\fi
% Use upquote if available, for straight quotes in verbatim environments
\IfFileExists{upquote.sty}{\usepackage{upquote}}{}
\IfFileExists{microtype.sty}{% use microtype if available
  \usepackage[]{microtype}
  \UseMicrotypeSet[protrusion]{basicmath} % disable protrusion for tt fonts
}{}
\makeatletter
\@ifundefined{KOMAClassName}{% if non-KOMA class
  \IfFileExists{parskip.sty}{%
    \usepackage{parskip}
  }{% else
    \setlength{\parindent}{0pt}
    \setlength{\parskip}{6pt plus 2pt minus 1pt}}
}{% if KOMA class
  \KOMAoptions{parskip=half}}
\makeatother
\usepackage{xcolor}
\usepackage[margin=2.1cm]{geometry}
\usepackage{longtable,booktabs,array}
\usepackage{calc} % for calculating minipage widths
% Correct order of tables after \paragraph or \subparagraph
\usepackage{etoolbox}
\makeatletter
\patchcmd\longtable{\par}{\if@noskipsec\mbox{}\fi\par}{}{}
\makeatother
% Allow footnotes in longtable head/foot
\IfFileExists{footnotehyper.sty}{\usepackage{footnotehyper}}{\usepackage{footnote}}
\makesavenoteenv{longtable}
\setlength{\emergencystretch}{3em} % prevent overfull lines
\providecommand{\tightlist}{%
  \setlength{\itemsep}{0pt}\setlength{\parskip}{0pt}}
\setcounter{secnumdepth}{5}
\ifLuaTeX
  \usepackage{selnolig}  % disable illegal ligatures
\fi
\usepackage{bookmark}
\IfFileExists{xurl.sty}{\usepackage{xurl}}{} % add URL line breaks if available
\urlstyle{same}
\hypersetup{
  colorlinks=true,
  linkcolor={Maroon},
  filecolor={Maroon},
  citecolor={Blue},
  urlcolor={Blue},
  pdfcreator={LaTeX via pandoc}}

\author{}
\date{}

\begin{document}

{
\hypersetup{linkcolor=}
\setcounter{tocdepth}{3}
\tableofcontents
}
\section{ThreatForge Decision regression after model-layering Revision
7}\label{threatforge-decision-regression-after-model-layering-revision-7}

\begin{quote}
\textbf{Non-canonical research checkpoint.}

Baseline of the study: Decision model Revision 7 after separating L1
conceptual semantics, L2 representation contracts, L3 realization
principles and L4 tool support/conformance. Construction corpus only:
the sixteen revised candidate ADRs historically associated with
ThreatForge MR-0001 and MR-0002. ThreatForge MR-0003, MR-0004 and
MR-0005 remain untouched holdouts.
\end{quote}

\subsection{Purpose}\label{purpose}

This checkpoint re-runs the construction corpus after discovering that
\texttt{R-CLOSED-04} had been used at the wrong level. Executable
projections remain a useful L3 realization principle, but they are not
an L1 semantic invariant inherited by every governed project. Therefore
a ThreatForge ADR is not redundant merely because it chooses a concrete
way to realize that principle.

Each candidate is reviewed through four questions:

\begin{enumerate}
\def\labelenumi{\arabic{enumi}.}
\tightlist
\item
  Is there a significant choice rather than a consequence, requirement
  or explanation?
\item
  Who owns the choice semantically: DDTA method/metamodel, canonical
  representation, ThreatForge product/tool, or ThreatForge project
  operations?
\item
  Is the historical parent Macro Requirement the natural semantic owner?
\item
  Does the candidate contribute a non-redundant choice without using
  L2/L3 principles as if they were L1 semantics?
\end{enumerate}

\subsection{Regression table}\label{regression-table}

\small

\begin{longtable}[]{@{}
  >{\raggedright\arraybackslash}p{(\columnwidth - 6\tabcolsep) * \real{0.2500}}
  >{\raggedright\arraybackslash}p{(\columnwidth - 6\tabcolsep) * \real{0.2500}}
  >{\raggedright\arraybackslash}p{(\columnwidth - 6\tabcolsep) * \real{0.2500}}
  >{\raggedright\arraybackslash}p{(\columnwidth - 6\tabcolsep) * \real{0.2500}}@{}}
\toprule\noalign{}
\begin{minipage}[b]{\linewidth}\raggedright
Historical slot
\end{minipage} & \begin{minipage}[b]{\linewidth}\raggedright
Revised candidate
\end{minipage} & \begin{minipage}[b]{\linewidth}\raggedright
Revision 7 result
\end{minipage} & \begin{minipage}[b]{\linewidth}\raggedright
Rationale
\end{minipage} \\
\midrule\noalign{}
\endhead
\bottomrule\noalign{}
\endlastfoot
\texttt{MR-0001\ /\ ADR-0001} & Diátaxis come classificazione comune
della funzione documentale & \textbf{RELOCATE - Decision DDTA method} &
Scelta reale (Diátaxis vs classificazioni locali), ma la candidata
attribuisce la scelta al metodo DDTA. Non è una Decision figlia naturale
del progetto/tool ThreatForge. \\
\texttt{MR-0001\ /\ ADR-0002} & Vocabolario governato come fonte unica
del linguaggio condiviso & \textbf{RELOCATE - DDTA/cross-document
candidate} & Scelta reale (significati governati vs convenzioni locali),
ma l'owner è DDTA/cross-document governance. Resta aperto se il
controlled vocabulary sia universale, opzionale o profile-specific. \\
\texttt{MR-0001\ /\ ADR-0003} & Supporto ThreatForge alla tracciabilità
della realizzazione governata & \textbf{REWRITE + REPARENT candidate} &
Tolto l'uso improprio di R-CLOSED-04, il tema può contenere una vera
scelta ThreatForge; la rewrite corrente però dichiara soprattutto
supporto senza scegliere abbastanza. Il concern è più naturale sotto
MR-0002. \\
\texttt{MR-0001\ /\ ADR-0004} & Risoluzione contestuale dei valori
controllati da fonti governate & \textbf{FITS} & Scelta ThreatForge
autonoma: risolvere i valori nel contesto posseduto dalle fonti
governate invece di codificare enum/global interpretations nei consumer.
L3 non prescrive semanticamente questa scelta. \\
\texttt{MR-0001\ /\ ADR-0006} & Handoff governato per la continuità
dello sviluppo ThreatForge & \textbf{FITS - parent unresolved} & Scelta
reale fra continuità informale e handoff governato ricostruibile. È una
ThreatForge project-operation; il parent MR corrente non va forzato e
deve essere riesaminato. \\
\texttt{MR-0001\ /\ ADR-0007} & Contratti eseguibili dei modelli
documentali consumati da ThreatForge & \textbf{FITS} & Scelta
architetturale reale: consumer che ricodificano separatamente
struttura/regole vs una proiezione/contratto eseguibile comune.
R-CLOSED-04 è L3 e non rende questa ADR semanticamente ridondante. \\
\texttt{MR-0001\ /\ ADR-0008} & Rappresentazione Doc-as-Code canonica
dei riferimenti governati & \textbf{RELOCATE - L2 representation
candidate} & La grammatica {[}id{]} title è una scelta reale di
rappresentazione, non una Decision del tool. Va rivalutata nella futura
fase identity/reference/mapping. \\
\texttt{MR-0001\ /\ ADR-0009} & Supporto ThreatForge a obblighi di
sicurezza governati derivati dall'analisi & \textbf{NOT STANDALONE
DECISION candidate} & Non viene eliminata da L3, ma H-CLOSED-02 resta
decisivo: nella rewrite corrente supportare Security tramite contratti
sembra coverage/conformance di ADR-0007 + ADR-0010, non una nuova scelta
autonoma. \\
\texttt{MR-0001\ /\ ADR-0010} & Consumer ThreatForge estensibili
dall'inventario governato dei modelli & \textbf{FITS} & Scelta autonoma:
inventario/dispatch derivato da catalogo/proiezione governata vs liste
locali, cardinalità fisse o discovery indipendente. Distinta dal modo in
cui i contratti sono consumati. \\
\texttt{MR-0002\ /\ ADR-0001} & Separazione tra capacità ThreatForge,
orchestrazione applicativa e adapter & \textbf{FITS - minor cleanup} &
Scelta architetturale autonoma. La Decision deve possedere il confine
capability/application/adapter senza anticipare quale adapter concreto
venga scelto; VS Code resta ADR-0005. \\
\texttt{MR-0002\ /\ ADR-0002} & Pubblicazione governata del repository
ThreatForge & \textbf{FITS} & Scelta di project operation: Git manuale
vs operazione governata con precondizioni. Il parent MR-0002 è
naturale. \\
\texttt{MR-0002\ /\ ADR-0003} & Supporto ThreatForge all'avvio governato
della realizzazione & \textbf{FITS} & Scelta di prodotto: scaffold
governato e collegato al contesto vs creazione tecnica scollegata.
Stati/lifecycle formali restano downstream. \\
\texttt{MR-0002\ /\ ADR-0004} & Authoring ThreatForge guidato dai
contratti governati del modello & \textbf{FITS} & Scelta capability
autonoma: sola validazione a posteriori vs assistenza proattiva derivata
dai contratti durante authoring. ADR-0007 può fornire la fonte senza
imporre questa capability. \\
\texttt{MR-0002\ /\ ADR-0005} & VS Code task come adapter locale sottile
di ThreatForge & \textbf{FITS} & Scelta concreta di prima integrazione:
task sottili vs custom extension/altro adapter. Distinta dal confine
architetturale generale di ADR-0001. \\
\texttt{MR-0002\ /\ ADR-0006} & Core ThreatForge condiviso per
assistenza Markdown contestuale & \textbf{FITS} & Scelta autonoma su
dove vive l'analisi live: core editor-independent vs regole replicate
negli editor. Distinta da authoring guidato e scelta dei task. \\
\texttt{MR-0002\ /\ ADR-0007} & Architettura applicativa ThreatForge
indipendente dai delivery adapter & \textbf{FITS} & Scelta
architetturale più specifica di ADR-0001: application services +
ports/adapters + composition boundary. Può cambiare indipendentemente
dal macro-confine capability/app/adapter. \\
\end{longtable}

\normalsize

\subsection{Cross-case findings}\label{cross-case-findings}

\subsubsection{F1 - Realization principles do not erase tool
Decisions}\label{f1---realization-principles-do-not-erase-tool-decisions}

\texttt{R-CLOSED-04} can motivate ThreatForge architecture but cannot be
used as an L1 inheritance rule. In particular, choosing shared
executable contracts rather than independently coded consumer rules
remains a real ThreatForge Decision.

\subsubsection{F2 - Ownership-correct is not
sufficient}\label{f2---ownership-correct-is-not-sufficient}

A rewrite that merely says `ThreatForge supports governed semantics
without owning them' may still fail the Decision-contribution test. The
corrected candidate must identify a genuine alternative in the behavior
or architecture of ThreatForge.

\subsubsection{F3 - Re-parent after ownership
correction}\label{f3---re-parent-after-ownership-correction}

When a rewrite changes the semantic owner or narrows the problem, its
parent relation must be re-evaluated. Historical parentage is evidence,
not an immutable truth.

\subsubsection{F4 - Support/conformance is not automatically a
Decision}\label{f4---supportconformance-is-not-automatically-a-decision}

Supporting one concrete model extension through already-chosen generic
contracts may be L4 coverage rather than another Decision.
MR-0001/ADR-0009 is the primary construction example.

\subsubsection{F5 - The construction corpus remains evidence, not
migration}\label{f5---the-construction-corpus-remains-evidence-not-migration}

This checkpoint does not modify historical ThreatForge ADRs. It
classifies revised candidates so that the next rewrite can test whether
the Decision metamodel produces a smaller, ownership-correct and
non-redundant set.

\subsection{Disposition summary}\label{disposition-summary}

\begin{itemize}
\tightlist
\item
  Retain and rewrite/clean as ThreatForge Decisions: MR-0001/ADR-0003,
  ADR-0004, ADR-0007, ADR-0010 and MR-0002/ADR-0001..ADR-0007.
\item
  Retain as a meaningful ThreatForge Decision but keep parent
  unresolved: MR-0001/ADR-0006 handoff.
\item
  Relocate outside the ThreatForge Decision hierarchy: MR-0001/ADR-0001
  Diátaxis, ADR-0002 vocabulary, ADR-0008 reference representation.
\item
  Remove as standalone Decision candidate and use as L4
  support/conformance case: MR-0001/ADR-0009 Security support.
\end{itemize}

\subsection{Next step}\label{next-step}

Write a corrected candidate corpus using
\texttt{Context\ +\ Decision\ +\ Consequences} for actual Decisions,
with study-only annotations for semantic owner, proposed parent and DDTA
contract supported. These annotations are not new canonical fields of
the Decision metamodel.

\end{document}
